
We opened with a fundamental question: when code generation models achieve 70--80\% on standard benchmarks but drop to 30--40\% on hidden tests (HumanEval $\rightarrow$ HumanEval+), does execution feedback truly align with human intent? Our systematic investigation of feedback orthogonality across task types provides the answer: alignment depends on specification completeness. For better-specified tasks (competitive programming), execution moderately proxies intent ($\rho$=0.68); for underspecified realistic tasks (software engineering), the proxy weakens dramatically ($\rho$=0.35 predicted, 2.$29\times$ variance). Specification completeness drives this gap through test-intent coverage: underspecified tasks miss 2.$00\times$ the intent dimensions (67\% versus 33\%, p<0.0001), leaving execution feedback blind to critical dimensions humans evaluate.

Supervised AI feedback offers a task-independent alternative. CodeBERT trained on human annotations achieves strong intent alignment ($\rho$=0.85, +75\% versus zero-shot conflating supervision and architecture advantage) independent of specification completeness, demonstrating that direct supervision on human judgment patterns bypasses task-dependency. This parallels InstructGPT's RLHF for text generation: supervised learning captures the full intent space (all dimensions) rather than just the tested subset (correctness/edge cases).

Our three core contributions advance code generation alignment:

\begin{enumerate}
\item \textbf{First systematic feedback orthogonality mapping}: Pairwise correlations (execution/AI/human) across task types (competitive/realistic) reveal task-dependent structure (ANOVA F=2226.34, p<0.0001). Prior work studied modalities in isolation (CodeRL execution-only, RLAIF AI-human for text); we compare all three for code across specification spectrum.
\end{enumerate}

\begin{enumerate}
\item \textbf{Specification completeness mechanism validation}: Qualitative dimension analysis (2.$00\times$ missed dimension gap, $\chi^{2}=53.33$ p<0.0001) explains why execution feedback fails for realistic tasks: tests capture only approximately 33\% of intent dimensions (functional correctness) while missing approximately 67\% (readability, maintainability, efficiency, security). Causal chain validated: specification completeness $\rightarrow$ test coverage $\rightarrow$ execution-human correlation.
\end{enumerate}

\begin{enumerate}
\item \textbf{Supervised AI path extending supervised learning paradigm}: We extend Li et al. (2022) supervised learning approach from code review to alignment feedback, showing that supervised CodeBERT achieves $\rho$=0.85 strong alignment, providing viable alternative to execution-only approaches. Demonstrates InstructGPT analogy for code---supervised training bypasses specification dependency through direct intent modeling.
\end{enumerate}

These findings challenge the execution-only alignment paradigm (CodeRL effective for HumanEval but task-specific) and enable practical alternatives. The path forward is adaptive: use task type to route feedback modalities where each provides strongest signal. Competitive tasks (complete specifications) can trust execution for moderate alignment at near-zero cost. Realistic tasks (underspecified requirements) require supervised AI feedback for strong alignment without human-in-the-loop during deployment.

\textbf{Future vision}: Immediate extensions include task type classifiers (predict competitive/realistic from problem text $\rightarrow$ automatic routing) and dimension-specific AI models (train separate models for correctness/readability/efficiency $\rightarrow$ ensemble voting for $\rho$>0.9 alignment). Medium-term, adaptive weighting systems can combine execution's runtime guarantees (catches functional errors) with AI's intent understanding (evaluates quality dimensions) through learned weighting functions. Long-term, multi-modal alignment pipelines integrating execution, supervised AI, and selective human feedback could achieve comprehensive coverage: execution for runtime correctness, AI for quality assessment, human validation for edge cases.

For code generation alignment, one size does not fit all. By quantifying when execution feedback suffices (better-specified tasks) and when it fails (underspecified tasks), while demonstrating task-independent supervised AI alternatives, we enable alignment strategies that adapt to specification completeness rather than assuming execution universally proxies intent. The 67\% of intent dimensions tests miss for realistic tasks---readability, maintainability, efficiency, security---are precisely the dimensions developers care about when code moves from benchmarks to production. Our work provides both understanding (specification mechanism) and tools (supervised AI path) to align on what matters.
